% !TEX root = ../main.tex

\section{Trabajo relacionado}
\label{sec:related}

\sculpt se inserta en una conversaci\'on que ya tiene varias d\'ecadas
sobre c\'omo abrir la programaci\'on a p\'ublicos no especializados.
Esta ha producido tres familias de lenguajes que conviene
revisar por separado, ya que cada una resuelve parcialmente el
problema que \sculpt plantea. Se nombran a estas familias
\emph{lenguajes visuales y de bloques}, \emph{lenguajes tangibles para
la ense\~nanza de la computaci\'on} y \emph{computaci\'on como pr\'actica
art\'istica} (\emph{creative computing}). Al cierre de la secci\'on
se posiciona a \sculpt frente a las tres.

\subsection{Lenguajes visuales y de bloques}
\label{sec:related:visual}

La l\'inea de los lenguajes basados en bloques arranca con
\emph{LogoBlocks}~\cite{begel96logoblocks} en 1996 y se consolida con
\emph{Scratch}~\cite{resnick+09scratch}, desarrollado en el
\emph{Lifelong Kindergarten Group} del MIT Media Lab. Scratch reemplaza
la sintaxis textual por bloques arrastrables cuya forma define qu\'e
encaja con qu\'e. Esto abre una exploración interesante sobre la sintaxis.
La sintaxis en este tipo de lenguajes adquiere una dimenison geometrica.
Los conjuntos de reglas que gobiernan la construcci\'on de programas son 
ahora reglas fisicas. Esa idea ha sido extendida por \emph{Blockly}~\cite{fraser15blockly},
una biblioteca de Google que genera c\'odigo en lenguajes objetivo
(JavaScript, Python, Dart) a partir de programas en bloques, y por
\emph{Snap!}~\cite{harvey13snap}, una reimplementaci\'on de Scratch
con primer-orden de funciones y bloques personalizables, orientada a la
educaci\'on universitaria.

Estas propuestas pretenden reducir dr\'asticamente la fricci\'on sint\'actica
para los principiantes y son hoy la herramienta de elecci\'on en la
ense\~nanza temprana de la computaci\'on. Sus limitaciones, sin embargo,
son congruentes con el diagn\'ostico de la problematica. Los
bloques viven en una pantalla, dependen de un teclado (usualmente) y un \emph{mouse},
y, por dise\~no, no admiten libertad est\'etica significativa
m\'as all\'a de cambiar colores de fondo. La interfaz se ha simplificado, la computación mejorado,
pero su soporte material no ha cambiado.

\subsection{Lenguajes tangibles para la ense\~nanza}
\label{sec:related:tangible}

Una segunda familia traslada la programaci\'on al mundo f\'isico.
\emph{AlgoBlock}~\cite{suzuki95algoblock} fue uno de los primeros
sistemas en proponer bloques f\'isicos conectables como medio para
construir programas, principalmente con prop\'ositos colaborativos
entre ni\~nos. \emph{Tern}~\cite{horn07tern} extendi\'o la idea con piezas
de madera reconocidas por visi\'on por computador, sin necesidad de
electr\'onica embebida. M\'as recientemente, \emph{KIBO}~\cite{sullivan15kibo}
combina bloques de madera con c\'odigos de barras leidos por un robot
m\'ovil, y \emph{Cubetto}~\cite{cubetto}, comercializado en m\'ultiples
mercados, usa una cuadricula y fichas magn\'eticas para programar el
desplazamiento de un peque\~no robot. \emph{Osmo Coding}~\cite{osmo}
sigue una l\'inea h\'ibrida, las piezas son f\'isicas pero la ejecuci\'on
ocurre en una tableta que las observa por la c\'amara.

Estos sistemas demuestran ampliamente que la tangibilidad reduce
barreras cognitivas y motivacionales en p\'ublicos estudiantiles, en
linea con literatura previa sobre \emph{tangible user 
interfaces}~\cite{ishii97tangiblebits,marshall07tangible}. Comparten, 
no obstante, tres limitaciones recurrentes. Primero, todos est\'an 
dise\~nados como \emph{herramientas educativas}: su poder expresivo es 
intencionalmente acotado y aplicado a un dominio específico, usualmente robotica. 
Esto brinda una experiencia de programaci\'on m\'as concreta, manipulativa y estumulante 
para los estudiantes, pero no dejan de ser entonces lenguajes de dominio espec\'ifico, o en otras 
palabras, herramientas didacticas programables. Segundo, las piezas son cerradas. 
Vienen en una forma, un color y un tama\~no preestablecidos por el 
fabricante, y el usuario no puede siempre modificarlas sin destruir o alterar su 
funcionalidad. Esta libertad, tanto est\'etica como funcional, por dise\~no esta ausente.
Este af\'an por controlar la experiencia del usuario es comprensible en un contexto educativo, 
sin embargo, limita la expresividad y la cantidad de usuarios que pueden usar estos sistemas.
Libertad de modificacion con las piezas es relevante para usuarios con necesidades de
accesibilidad, por ejemplo, que pueden necesitar piezas de un tama\~no, material o morfolog\'ia
particular para manipularlas. Tercero, la mayor\'ia de estos sistemas vienen con una etiqueta de
precio que no todas las instituciones educativas o individuos pueden pagar. Adicional a necesitar
hardware (y a veces software) especializado, el costo de las piezas, su mantenimiento y acceso puede ser un
obst\'aculo para su adopci\'on. Esto es especialmente relevante en contextos de bajos recursos,
donde la accesibilidad a herramientas educativas resulta crucial para cerrar brechas de aprendizaje.

\subsection{Computaci\'on como pr\'actica art\'istica}
\label{sec:related:art}

Una tercera tradici\'on aborda el problema desde un \'angulo distinto.
Mantener la interfaz textual de la programaci\'on pero dirigida
expl\'icitamente a la producci\'on art\'istica. \emph{Processing}, de
Reas y Fry~\cite{reas07processing}, ofrece un dialecto de Java pensado
para artistas visuales y dise\~nadores, y populariz\'o la idea de que
escribir c\'odigo es una forma leg\'itima de pr\'actica art\'istica.
\emph{Sonic Pi}, de Aaron y Blackwell~\cite{aaron16sonicpi}, construye
sobre el motor de s\'intesis SuperCollider un superconjunto de Ruby
orientado a \emph{live coding} musical, originalmente concebido como
herramienta educativa para escuelas en Inglaterra. La pr\'actica del
\emph{live coding}, tanto musical como visual, ha sido sistematizada
como movimiento art\'istico-computacional propio~\cite{collins03livecoding},
con festivales y comunidades dedicadas (\eg \emph{Algorave}).

Estas iniciativas han demostrado que la programaci\'on \emph{puede} ser
arte y que existe audiencia entrenada y entusiasta para esa visi\'on.
Sin embargo, comparten un l\'imite estructural con la primera familia.
No abandonan el paradigma de la m\'aquina de escribir. El artista que
hace \emph{live coding} sigue tecleando sobre una pantalla en vivo (y estresandose porque el c\'odigo no compila)
art\'istica final es sonido, video o imagen, no el cuerpo f\'isico del
programa.

\subsection{Posicionamiento de \sculpt}
\label{sec:related:positioning}

\sculpt se sit\'ua en la intersecci\'on poco transitada de las tres
familias anteriores. Comparte con los lenguajes de bloques la idea de
que la sintaxis puede ser una propiedad geom\'etrica de las piezas
esencial para la construcci\'on de programas. La intuitividad natural de "encajar" piezas
es un recurso poderoso para asimilar conceptos sint\'acticos complejos y los aterriza a una
experiencia y logica de construcci\'on concreta y entendible. Encaja o no encaja.
Comparte con los lenguajes tangibles la apuesta por sacar la programaci\'on de la 
pantalla y llevarla al espacio tridimensional.
Habilitarle al mundo a programar implica necesariamente abrir la programaci\'on a formas
de acceso distintas a la tradicional. Se pretende aprovechar la naturaleza manipulativa y abierta
de las fichas para abrir la programaci\'on a usuarios con distintas capacidades, recursos
y necesidades de accesibilidad.
Finalmente, comparte con la tradici\'on \emph{creative computing} la convicci\'on
de que la pr\'actica art\'istica es un veh\'iculo leg\'itimo para la
computaci\'on, no un adorno posterior. La tabla \ref{tab:related} resume
las diferencias.

\begin{table}[h]
\centering
\small
\caption{Comparaci\'on de \sculpt con familias relacionadas.}
\label{tab:related}
\begin{tabular}{@{}lcccc@{}}
\toprule
                          & Bloques/visuales & Tangibles edu. & Creative comp. & \sculpt \\
\midrule
Soporte f\'isico          & No  & S\'i & No  & S\'i \\
Turing-completo           & Parcial & No & S\'i & S\'i \\
Libertad est\'etica       & Baja & Baja & Media & Muy Alta \\
P\'ublico objetivo        & Estudiantes/Novatos & Ni\~nos & Artistas & Mixto \\
\bottomrule
\end{tabular}
\end{table}

Hasta donde se sabe, \sculpt es el primer lenguaje que combina
\emph{tangibilidad f\'isica}, \emph{Turing-completitud demostrada
formalmente} y \emph{libertad est\'etica radical} en un mismo dise\~no.
La radicalidad de la apertura est\'etica merece subrayarse: en \sculpt
el usuario puede modificar libremente forma, tama\~no, color, material
y textura de las piezas; lo \'unico restringido es la topolog\'ia de
las conexiones. Una pieza puede ser un cubo PLA est\'andar o una figura
animal tallada en madera, y el programa funcionar\'a igual mientras los
puntos de acople respeten la especificaci\'on.

\endinput
